\documentclass[a4paper,10pt,uplatex,dvipdfmx,twocolumn]{jsarticle}
\usepackage[a4paper,top=24mm,bottom=25mm,left=15mm,right=15mm,nohead,nofoot]{geometry}
\usepackage{amsmath,amssymb,graphicx,tikz,pgfplots,float,titlesec,url}
\usepackage{cite}
\usepackage{etoolbox}
\apptocmd{\thebibliography}{\setlength{\itemsep}{0pt}\setlength{\parsep}{0pt}}{}{}
\usetikzlibrary{arrows.meta,positioning}
\pgfplotsset{compat=1.18}
\setlength{\columnsep}{8mm}\setlength{\emergencystretch}{2em}\hbadness=3500
\setlength{\parskip}{0pt}
\setlength{\intextsep}{3pt}
\setlength{\textfloatsep}{6pt plus 2pt minus 2pt}
\setlength{\abovecaptionskip}{2pt}
\setlength{\belowcaptionskip}{0pt}
\titleformat{\section}{\normalsize\bfseries}{\thesection.}{.5em}{}
\titlespacing*{\section}{0pt}{0pt}{0pt}
\pagestyle{empty}
\makeatletter
\renewcommand\@biblabel[1]{#1)}
\renewcommand\citeleft{}
\renewcommand\citeright{)}
\renewcommand\citepunct{,}
\renewcommand\citemid{,\ }
\makeatother
\AtBeginDvi{\special{pdf:mapline uprml-h unicode ipaexm.ttf}\special{pdf:mapline uprml-hq unicode ipaexm.ttf}\special{pdf:mapline upgbm-h unicode ipaexg.ttf}\special{pdf:mapline upgbm-hq unicode ipaexg.ttf}\special{pdf:tounicode UTF8-UTF16}}
\begin{document}
\setlength{\abovedisplayskip}{0pt}\setlength{\abovedisplayshortskip}{0pt}\setlength{\belowdisplayshortskip}{0pt}
\setlength{\belowdisplayskip}{0pt}
\fontsize{10}{14.6}\selectfont
\twocolumn[\noindent\hspace*{36mm}{\Large\bfseries 収集画像からの部分画像生成とLoRAによるPTP包装識別}\par
\vspace{14.6pt}
\noindent\hfill 電子商取引研究室　呉竹 櫂人，森山 真光\hspace{1zw}\par
\vspace{14.6pt}]
\section{序論}
薬剤の取り違えは患者に重い健康被害を与えるおそれがあり，PTP包装の画像から薬剤を識別する支援手法が研究されている\cite{who2017,han2021}。
しかし，従来研究では，1薬剤あたり表裏各72枚，計144枚の画像を撮影して学習に用いている\cite{han2021}。
本研究の目的は，各薬剤の表裏各1枚の収集画像だけから，実写画像の薬剤名を識別することである。
そのために，収集画像から部分画像と拡張画像を作り，視覚言語モデルをLoRAで追加学習する方法を提案する。

\section{部分画像とLoRAによる薬剤名識別}
図\ref{fig:flow}に，提案手法の学習と識別の流れを示す。
学習では，表裏各1枚の収集画像から，シート全体と，1錠区画・2分割・4分割の部分画像を作る。
各画像に回転12種と明るさ5種の変換を別々に施し，1枚から計17枚の拡張画像を作る。
モデルには画像と指示文「この錠剤は何という薬ですか？薬剤名だけで答えてください。」を入力し，正解の薬剤名を出力するように学習する。
識別時も同じ指示文を使い，薬剤名の候補は与えない。
図\ref{fig:maintate}に，メインテート錠2.5mgの収集画像，部分画像，実写画像の例を示す。

\begin{figure}[H]
\centering
\begin{tikzpicture}[x=1mm,y=0.85mm,font=\fontsize{8}{9}\selectfont,
 act/.style={draw,rounded corners=1.5mm,align=center,minimum height=5mm,inner sep=1.5pt},
 obj/.style={draw,align=center,minimum height=5mm,inner sep=1.5pt},
 arr/.style={-{Stealth[length=1.5mm]},line width=.45pt}]
\draw[gray] (0,0) rectangle (86,-33);
\draw[gray] (43,0)--(43,-33);
\node at (21.5,-2.5) {学習};
\node at (64.5,-2.5) {識別};
\fill (21.5,-6) circle (1.2mm); \fill (64.5,-6) circle (1.2mm);
\node[act] (a) at (21.5,-11.5) {部分画像を生成する};
\node[act] (b) at (21.5,-18.5) {回転・明るさで拡張する};
\node[act] (c) at (21.5,-25.5) {LoRAを学習する};
\draw (21.5,-30.5) circle (1.5mm); \fill (21.5,-30.5) circle (.9mm);
\node[act] (e0) at (64.5,-11.5) {実写画像を入力する};
\node[act] (e) at (64.5,-18.5) {薬剤名を生成する};
\node[act] (f) at (64.5,-25.5) {正解名と照合する};
\draw (64.5,-30.5) circle (1.5mm); \fill (64.5,-30.5) circle (.9mm);
\node[obj] (o) at (43,-25.5) {学習済み\\LoRA};
\draw[arr] (21.5,-7.2)--(a); \draw[arr] (64.5,-7.2)--(e0);
\draw[arr] (a)--(b); \draw[arr] (b)--(c);
\draw[arr] (e0)--(e); \draw[arr] (e)--(f);
\draw[arr] (c)--(o); \draw[arr] (o.north) |- (e.west);
\draw[arr] (c)--(21.5,-29); \draw[arr] (f)--(64.5,-29);
\end{tikzpicture}
\caption{提案手法の学習と識別の流れ（UMLアクティビティ図）}
\label{fig:flow}
\end{figure}
\begin{figure}[H]
\centering
\includegraphics[width=\linewidth]{fig_maintate.png}
\caption{収集画像から作る部分画像と実写画像の例}
\label{fig:maintate}
\end{figure}

LoRAは，元の重み$W_0$を固定し，低ランク行列の積を更新として加える手法である\cite{hu2022}。
学習するのは，式\eqref{eq:lora}の行列$A$と$B$だけである。
\begin{equation}\label{eq:lora}
W=W_0+\frac{\alpha}{r}BA
\end{equation}
ここで，$r$は$A$と$B$のランク，$\alpha$は更新の大きさを調整する係数である。
識別結果は，表記をそろえた回答に正解の薬剤名が含まれていれば正解とし，正解率で評価する。

\section{実験結果と考察}
表\ref{tab:settings}に，実験条件を示す。
学習には，Webで収集した41薬剤の表裏各1枚，計82枚から作った1,224枚を用いた。
評価には，近畿大学病院薬剤部で撮影した41薬剤の実写画像287枚を用いた。
評価画像は，学習には使わなかった。
比較として，追加学習していないGemini 3.8 Flashを2条件で評価した。
短文条件では提案手法と同じ指示文を与え，候補一覧条件では41薬剤の一覧から1つを選ばせた。
\begin{table}[t]
\centering
\caption{実験条件}
\label{tab:settings}
\fontsize{8}{9}\selectfont
\setlength{\tabcolsep}{2pt}
\begin{tabular}{@{}p{1.7cm}p{6.6cm}@{}}
\hline
基盤モデル & Qwen3.5-4B \\
LoRA & $r=16$，$\alpha=128$ \\
学習 & 総提示107,525枚，学習率$2\times10^{-4}$，BF16 \\
計算環境 & NVIDIA H100 1基，PyTorch 2.11.0，PEFT 0.19.1 \\
比較モデル & Gemini 3.8 Flash（temperature 0） \\
\hline
\end{tabular}
\end{table}
図\ref{fig:accuracy}に，各条件の正解率を示す。
提案手法は，学習の乱数を変えた3回で，287枚中286枚，287枚，284枚に正解した。
平均正解率は99.54\%，標準偏差は0.53ポイントであった。
Geminiは，短文条件で238枚（82.93\%），候補一覧条件で283枚（98.61\%）に正解した。
\begin{figure}[H]
\centering
\begin{tikzpicture}
\begin{axis}[
width=7.8cm,height=25mm,ybar,bar width=8mm,
ymin=0,ymax=125,ytick={0,50,100},
ylabel={正解率（\%）},
symbolic x coords={Q,G1,G2},xtick={Q,G1,G2},
xticklabels={提案手法,Gemini短文,Gemini候補一覧},
enlarge x limits=.2,
tick label style={font=\fontsize{8}{9}\selectfont},
label style={font=\fontsize{8}{9}\selectfont},
axis x line*=bottom,axis y line*=left,ymajorgrids,grid style={gray!20},
nodes near coords,point meta=y,
nodes near coords style={font=\fontsize{8}{9}\selectfont,/pgf/number format/fixed,/pgf/number format/precision=2,/pgf/number format/zerofill},
error bars/y dir=both,error bars/y explicit,
error bars/error bar style={black,line width=1pt},
error bars/error mark options={rotate=90,mark size=4pt,line width=1pt}]
\addplot[fill=black!65,draw=black,bar shift=0pt] coordinates {(Q,99.5354239257) +- (0,.5322387567)};
\addplot[fill=black!20,draw=black,bar shift=0pt] coordinates {(G1,82.9268292683) (G2,98.606271777)};
\end{axis}
\end{tikzpicture}
\caption{実写287枚での正解率（提案手法は3回の平均，誤差棒は標準偏差）}
\label{fig:accuracy}
\end{figure}
学習の各回について，Geminiとの差をMcNemar検定で調べた。
提案手法は，短文条件のGeminiを3回とも有意に上回った（$p<0.001$）。
候補一覧条件のGeminiとは，3回とも有意差がなかった（$p=0.375,\ 0.125,\ 1.000$）。
ただし，候補一覧条件は入力条件が異なるため，同じ性能とまでは判断できなかった。
なお，LoRAの設定は，評価画像で8条件を比べて平均正解率が最も高いものを選んだ。

\section{現状と今後の計画}
現在までに，41薬剤では，表裏各1枚の収集画像から実写画像を平均99.54\%で識別できた。
学習する薬剤を100種類に増やした場合，既知の41薬剤の実写287枚での平均正解率は96.75\%であった。
今後は，第1に，学習する薬剤を増やしても正解率を保てる方法と，モデルを小型にして正解率を上げる方法を検討する。
第2に，部分画像の内容による効果と，学習量が増えた効果を分けて確かめる。
第3に，新たに撮影した実写画像で評価し，条件選びに使っていない画像での正解率を確かめる。
あわせて，未学習の薬剤に対する回答の保留と，複数の薬剤が写った実写画像での評価にも取り組む。

{\fontsize{7}{8}\selectfont\raggedright
\bibliographystyle{junsrt}
\bibliography{261006}
}
\end{document}
